% 연습문제 모음 — 제4장
% 장은 이 파일의 이름으로 결정된다.

\begin{exo}[경로를 따라 미분형식 적분하기]{integrale-forme-differentielle}
\keywords{미분형식, 완전미분, 슈바르츠 판정법, 선적분}

두 미분형식
\[
\omega_1 = y\,\mathrm{d}x + x\,\mathrm{d}y,
\qquad
\omega_2 = y\,\mathrm{d}x,
\]
과 원점 $O(0,0)$에서 점 $M(1,1)$로 이어지는 세 경로를 생각하자.
$\gamma_1$은 수평선분 $O\to(1,0)$ 뒤에 수직선분 $(1,0)\to M$이 이어지는 경로이고,
$\gamma_2$는 수직선분 $O\to(0,1)$ 뒤에 수평선분 $(0,1)\to M$이 이어지는 경로이며,
$\gamma_3$은 $O$와 $M$을 직접 잇는 직선분이다.

\begin{enumerate}
\item $\omega_1 = \mathrm{d}f$를 만족하는 함수 $f$를 제시하여 $\omega_1$이 완전미분임을 보여라.
\item 각 선분을 매개화하여 $\int_{\gamma_1}\omega_1$과 $\int_{\gamma_2}\omega_1$을 계산하라. 계산 없이도 결과를 다시 구하라.
\item $\omega_2$에 슈바르츠 판정법을 적용하라. 무엇을 결론지을 수 있는가?
\item $\int_{\gamma_1}\omega_2$, $\int_{\gamma_2}\omega_2$, $\int_{\gamma_3}\omega_2$를 계산하고 논평하라.
\end{enumerate}

\begin{indication}
수평선분에서는 $\mathrm{d}y = 0$, 수직선분에서는 $\mathrm{d}x = 0$이므로
각 적분은 보통의 한 변수 적분으로 줄어든다.
\end{indication}

\begin{solution}
\textbf{문제 1.} $f(x,y) = xy$가 적합하다. 실제로 $\partial f/\partial x = y$이고 $\partial f/\partial y = x$이다.

\textbf{문제 2.} $t\mapsto\bigl(x(t),y(t)\bigr)$로 매개화된 곡선에 대해
\[
\int_\gamma\omega_1
= \int \left[y(t)x'(t)+x(t)y'(t)\right],\mathrm{d}t.
\]

경로 $\gamma_1$은 두 선분의 합이다. 첫 선분에서는 $x(t)=t$, $y(t)=0$이고
$t\in[0,1]$이므로 $\mathrm{d}x=\mathrm{d}t$, $\mathrm{d}y=0$이다.
둘째 선분에서는 $x(t)=1$, $y(t)=t$이고 역시 $t\in[0,1]$이므로
$\mathrm{d}x=0$, $\mathrm{d}y=\mathrm{d}t$이다. 따라서
\[
\begin{aligned}
I_1
= \int_{\gamma_1}\omega_1
&= \int_0^1\left(0\times 1+t\times 0\right)\,\mathrm{d}t
 + \int_0^1\left(t\times 0+1\times 1\right)\,\mathrm{d}t \\
&= 0+\int_0^1\mathrm{d}t = 1.
\end{aligned}
\]

$\gamma_2$의 첫 선분은 $x(t)=0$, $y(t)=t$, 둘째 선분은
$x(t)=t$, $y(t)=1$로 매개화하며 두 경우 모두 $t\in[0,1]$이다. 그러므로
\[
\begin{aligned}
I_2
= \int_{\gamma_2}\omega_1
&= \int_0^1\left(t\times 0+0\times 1\right)\,\mathrm{d}t
 + \int_0^1\left(1\times 1+t\times 0\right)\,\mathrm{d}t \\
&= 0+\int_0^1\mathrm{d}t = 1.
\end{aligned}
\]
두 값은 같다. 계산 없이도 예상할 수 있다. $\omega_1=\mathrm{d}f$이고
$f(x,y)=xy$이므로 적분은 끝점에만 의존하여
\[
\int_\gamma\omega_1=f(M)-f(O)=f(1,1)-f(0,0)=1.
\]

\textbf{문제 3.} $\omega_2=A(x,y)\,\mathrm{d}x+B(x,y)\,\mathrm{d}y$로 쓰자.
여기서 $A(x,y)=y$, $B(x,y)=0$이다. 이 형식이 완전미분이라면 슈바르츠 판정법에 따라
\[
\frac{\partial A}{\partial y}=\frac{\partial B}{\partial x}
\]
여야 한다. 그러나 $\partial A/\partial y=1$인 반면 $\partial B/\partial x=0$이다.
교차 편미분이 다르므로 $\omega_2$는 완전미분이 아니다. 따라서 그 적분은 $O$와 $M$
사이에서 택한 경로에 의존할 수 있으며, 다음 계산이 이를 확인한다.

\textbf{문제 4.} $\omega_2=y\,\mathrm{d}x$이므로 $y$가 0이 아닐 때의
$x$ 변화만 적분에 기여할 수 있다. 앞의 매개화를 사용하면 $\gamma_1$에서
\[
\begin{aligned}
J_1
=\int_{\gamma_1}\omega_2
&=\int_0^1 0\times 1\,\mathrm{d}t
  +\int_0^1 t\times 0\,\mathrm{d}t \\
&=0.
\end{aligned}
\]
$\gamma_2$의 첫 선분에서는 $x(t)=0$이므로 $\mathrm{d}x=0$이다.
둘째 선분에서는 $x(t)=t$, $y(t)=1$, $\mathrm{d}x=\mathrm{d}t$이므로
\[
\begin{aligned}
J_2
=\int_{\gamma_2}\omega_2
&=\int_0^1 t\times 0\,\mathrm{d}t
  +\int_0^1 1\times 1\,\mathrm{d}t \\
&=1.
\end{aligned}
\]
마지막으로 대각선분 $\gamma_3$은
\[
x(t)=t,\qquad y(t)=t,\qquad t\in[0,1]
\]
로 매개화된다. 따라서 $\mathrm{d}x=\mathrm{d}t$, $\mathrm{d}y=\mathrm{d}t$이고
\[
J_3
=\int_{\gamma_3}\omega_2
=\int_0^1 y(t)x'(t)\,\mathrm{d}t
=\int_0^1 t\,\mathrm{d}t
=\frac12.
\]
세 경로는 끝점이 같지만 $J_1=0$, $J_2=1$, $J_3=1/2$로 서로 다른 값을 준다.
이것이 완전미분이 아닌 미분형식의 특징적인 경로 의존성이다.
\end{solution}
\end{exo}

\begin{exo}[같은 내부 에너지 변화, 세 가지 전달 방식]{meme-delta-u-trois-transferts}
\keywords{열역학 제1법칙, 내부 에너지, 열, 역학적 일, 전기적 일, 열량 측정, 상태 함수}

액체와 열량계, 액체에 잠긴 작은 저항기가 거시적으로 정지한 닫힌계를 이룬다.
용기는 단단하며 계의 총 열용량은 일정하다고 가정한다:
\[
C=2{,}00\ \mathrm{kJ\,K^{-1}}.
\]
계를 $T_A=20{,}0\,^\circ\mathrm{C}$인 평형상태 $A$에서
$T_B=25{,}0\,^\circ\mathrm{C}$인 평형상태 $B$로 바꾸려 한다. 다음을 인정한다:
\[
U(B)-U(A)=C(T_B-T_A).
\]
거시적 운동 에너지와 위치 에너지의 변화는 0이다. 외부로의 모든 손실은 무시한다.

같은 초기상태 $A$에서 다음 세 절차를 각각 독립적으로 시행한다:
\begin{itemize}
\item[절차 1] 계에 아무 일도 공급하지 않고 고온 열원과 열접촉시킨다.
\item[절차 2] 벽을 단열한다. 질량 $m$인 추를 높이 $h=5{,}00$ m에서 떨어뜨려 날개가 액체를 젓게 한다. 마지막에 날개와 액체는 정지하며, 추가 잃은 모든 역학적 에너지는 계로 전달된다. $g=9{,}81\ \mathrm{m\,s^{-2}}$로 둔다.
\item[절차 3] 계는 열 $Q_3=4{,}00$ kJ를 받고, 나머지 에너지는 전기적으로 저항기에 공급된다. 계가 받는 전력은 일정하며 $\mathcal P=300$ W이다.
\end{itemize}

\begin{enumerate}
\item 세 절차에 공통인 내부 에너지 변화 $\Delta U=U(B)-U(A)$를 계산하라.
\item 첫 두 절차 각각에 대해 계가 받은 열 $Q$와 일 $W$를 구하라. 둘째 절차에서는 필요한 질량 $m$을 계산하라.
\item 셋째 절차에서 받은 전기적 일과 저항기에 전원을 공급한 시간을 계산하라.
\item 세 과정은 초기상태와 최종상태가 같다. 그런데도 $Q$와 $W$가 달라질 수 있는 이유를 설명하라. 최종상태 $B$에서 계는 ‘열’이나 ‘일’을 담고 있는가?
\end{enumerate}

\begin{indication}
은행가 부호 규약으로 $\Delta U=Q+W$를 적용하라. 날개의 경우 계가 받는 일은
추의 위치 에너지 감소량 $mgh$와 같다. 전선을 통해 경계를 지나는 전기 에너지는 전기적 일로 센다.
\end{indication}

\begin{solution}
\textbf{문제 1.} 온도 차이는 $T_B-T_A=5{,}0$ K이다. 따라서
\[
\Delta U=C(T_B-T_A)
=2{,}00\times5{,}0
=10{,}0\ \mathrm{kJ}.
\]
이 값은 두 평형상태 $A$와 $B$에만 의존한다.

\textbf{문제 2.} 첫 절차에서는 받은 일이 없으므로 $W_1=0$이다. 제1법칙에서
\[
Q_1=\Delta U=10{,}0\ \mathrm{kJ}.
\]
에너지가 계로 들어가므로 열은 양수이다.

둘째 절차에서는 벽이 단열되어 $Q_2=0$이다. 내부 에너지 변화는 모두 날개의 역학적 일에서 온다:
\[
W_2=\Delta U=10{,}0\ \mathrm{kJ}.
\]
$W_2=mgh$이므로
\[
m=\frac{W_2}{gh}
=\frac{10{,}0\times10^3}{9{,}81\times5{,}00}
\simeq 204\ \mathrm{kg}.
\]
이 큰 규모는 작은 온도 상승에도 이미 상당한 역학적 에너지가 대응함을 보여 준다.

\textbf{문제 3.} 제1법칙은
\[
W_3=\Delta U-Q_3=10{,}0-4{,}00=6{,}00\ \mathrm{kJ}
\]
를 요구한다. 전원이 계에 에너지를 주므로 이 일은 양수이다.
$W_3=\mathcal P\,\Delta t$에서
\[
\Delta t=\frac{W_3}{\mathcal P}
=\frac{6{,}00\times10^3}{300}
=20{,}0\ \mathrm{s}.
\]

\textbf{문제 4.} 세 경로는 각각
\[
(Q_1,W_1)=(10{,}0,0)\ \mathrm{kJ},
\qquad
(Q_2,W_2)=(0,10{,}0)\ \mathrm{kJ},
\]
그리고
\[
(Q_3,W_3)=(4{,}00,6{,}00)\ \mathrm{kJ}
\]
를 준다. 모든 경우 $Q+W$는 $10{,}0$ kJ이며 같은 $\Delta U$를 만든다.
내부 에너지는 상태 함수인 반면, 열과 일은 과정 중 에너지 전달을 나타낸다.
최종상태 $B$에서 계는 내부 에너지를 가지지만 ‘열’이나 ‘일’을 담지는 않는다.
\end{solution}
\end{exo}

\begin{exo}[일정한 외부 압력 아래의 압축]{compression-pression-exterieure-constante}
\seoready{true}
\keywords{압력 힘의 일, 일정한 외부 압력, 압축, 팽창, 진공 팽창}

기체를 일정한 외부 압력 $P_0$ 아래에서 부피 $V_A$로부터 $V_B<V_A$까지 압축한다.

\begin{enumerate}
\item 기체가 받은 일을 계산하라.
\item 기체가 같은 외부 압력 $P_0$에 맞서 $V_B$에서 $V_A$로 돌아간다. 받은 일을 계산하고 압축 때의 일과 비교하라.
\item 기체가 진공에서 $V_B$로부터 $V_A$까지 팽창하는 경우를 다시 다루어라. 얻은 세 부호를 해석하라.
\end{enumerate}

\begin{indication}
$W=-\int P_{\mathrm{ext}}\,\mathrm{d}V$를 사용하라. 과정이 준정적이지 않을 때도
적분해야 하는 것은 \emph{외부} 압력이다.
\end{indication}

\begin{solution}
\textbf{문제 1.} 외부 압력이 일정하므로
\[
W_{A\to B}=-P_0(V_B-V_A)=P_0(V_A-V_B)>0.
\]
압축할 때 기체는 일을 받는다.

\textbf{문제 2.} 되돌아가는 과정에서는
\[
W_{B\to A}=-P_0(V_A-V_B)<0.
\]
두 과정이 같은 일정한 외부 압력에 맞서 일어나므로 두 일은 서로 반대이다.

\textbf{문제 3.} 진공에서는 $P_{\mathrm{ext}}=0$이므로
\[
W_{B\to A}^{\mathrm{vide}}=0.
\]
따라서 팽창이 자동으로 음의 일을 뜻하지는 않는다. 실제로 외력을 밀어내야 한다.
\end{solution}
\end{exo}
